\scanpage{128}
% archive sheet number: 127
% source lines 1--10
$B_{\handwritten{\alpha}}$ sur les $A_{\handwritten{\alpha}}^{\handwritten{'}}$ qui se réduisent les unes en les autres \ldots), \deleted{\unsure{les $B$}} de plus chaque $B_{\handwritten{\alpha}}$ étant coplat sur $A_{\handwritten{\alpha}}$. La structure\inserted{\handwritten{$c$}} précédente revient à un isomorphisme de \uline{\deleted{bi}\inserted{co}modules} sur \deleted{$\unsure{B'|B|A'}\otimes_{\unsure{k'|k}}$} $A'$
\[
B\otimes_{k'}B\handwritten{\xleftarrow{\sim}}B\otimes_{k'}A'.\quad\deleted{\text{De}}
\tag*{(*)}
\]
Utilisant l'augmentation du facteur $A'$, on en conclut un homomorphisme
\[
B\otimes_{k'}B\handwritten{\longrightarrow}B
\tag*{(**)}
\]
qui sera un \uline{homomorphisme de comodules} sur $A'$ (NB utiliser la multiplication de $A'$ pour définir la structure de $A'$-comodule sur le produit tensoriel sur $k$ de deux tels), i.e. on a une loi de composition sur $B$ ``compatible avec la structure de $A^{\unsure{\mathrm{L}|'}}$-comodule''.
\margin{NB \unsure{Cette|Cettx} équivaut à dire que $B\handwritten{\to}B\otimes_{\unsure{k'|k}}A'$ est hom. d'alg.}
\margin{A préciser}
\margin{\handwritten{$\Vert$}Sauf erreur, la donnée de (*) \uline{équivaut à celle de} (**)}

% source lines 11--27
Si $\handwritten{\Delta}$ sur $U$ est unitaire i.e. $A$ est unitaire i.e. le $\otimes$ sur $X$ \deleted{\unsure{est adme}\illegible{1}} admet une unité bilatère compatible avec $F$, alors la compatibilité de $c$ avec les unités s'exprime, sauf erreur, par la condition que la loi de composition (**) est unitaire.\deleted{\unsure{(}} En effet, il faut écrire que l'on a un isomorphisme \deleted{$\unsure{L}\otimes_{\handwritten{U}}k$} $L\otimes_{U'}k'\handwritten{\simeq}k'$ satisfaisant certaines conditions, or cela revient à se donner une\deleted{élément de} augmentation de $L$ satisfaisant une certaine condition (savoir \deleted{être compatible avec}\inserted{être unité bilatère pour} la structure diagonale de $L$) ou encore un élément de $B$ qui es t une unité bilatère.\inserted{Ceci implique aussi que $B$ est \uline{fidèlement} plat sur $k'$.} De même, si $\handwritten{\Delta}$ est associatif, i.e.\inserted{si} on a pour $\otimes$ une \unsure{dxonnée|donnée} d'associatilité $\handwritten{\phi}$ compatible avec $F$, alors $c$ est compatible avec celle-ci sss $B$ est une algèbre associative. \deleted{Quand les deux conditions sont \unsure{vérifi|vérifié}} Enfin, quand $\handwritten{\Delta}$ est commutative i.e. $C$ avec $\handwritten{\psi}$ donnée de symétrie, \ldots\ alors $c$ est compatible avec cette dernière, et la donnée de symétrie habituelle sur $\unsure{\mathrm{Poojf}|\mathrm{Projf}}(k')$ sss la loi de composition de $B$ est \deleted{\unsure{asso}} commutative. \handwritten{[}(Mais alors, en termes du monoïde affine $\handwritten{\unsure{H|G|\mathscr{H}|\mathscr{G}}}$ défini par $A$ sur $k$, on voit que les structures de $B$ s'expriment en termes de\inserted{\handwritten{$P:$}} $\mathrm{Spec}\ B$ en disant que $\handwritten{H_{k'}}$ y opère à droite,\inserted{\handwritten{et $P$ est \unsure{form.|fibr.} princ. hom.}} et on voit de plus qu'après le changement
\margin{\handwritten{voir si $\unsure{C|G}$ un groupe i.e. si $\exists$ antipodisme.}}
